\chapter{Finite-stage carrier and base-change comparison}
\label{route-ch-finite-stage}

Return to the original curve \(C/K\).  Fix the separable closure
\(\Omega/K\) of Hypothesis H3, put \(C_\Omega=C\times_K\Omega\), and apply
Chapter~\ref{route-ch-sepclosed} to obtain the exact pair
\((J_\Omega,r_\Omega)\).  The finite-stage package is parameterized in the
Lean order \(\operatorname{Pic0FiniteStageGluePackage}\ C_\Omega\ K\): its
first curve argument is over the separably closed field \(\Omega\), while its
ambient field parameter is the original \(K\), with
\([\operatorname{Algebra}\ K\ \Omega]\) and
\([\operatorname{Algebra.IsAlgebraic}\ K\ \Omega]\).  Thus every displayed
\(P_0.N.1\) is a finite subextension of \(\Omega/K\), not a finite extension
of an unspecified field.

\section{A finite affine atlas}

\begin{theorem}[Finite-presentation atlas and simultaneous models]
  \label{route-thm-finite-atlas}
  \lean{AlgebraicGeometry.Pic0FiniteStageAtlas,
    AlgebraicGeometry.pic0FiniteStageAtlas,
    AlgebraicGeometry.finitePresentation_pic0FiniteStageChartRing,
    AlgebraicGeometry.exists_finSubext_pic0FiniteStageAtlas_chartRing_models}
  \leanok
  \source{stacks-project}
  \dcref{Tag 01ZM, especially parts (1)--(3)}
  \uses{route-thm-sepclosed-finiteness}

  Choose a finite set \(I\) of affine opens \(U_i\subset J_\Omega\) covering
  \(J_\Omega\).  Since the represented Picard group is separated, each
  ordered intersection \(U_{ij}:=U_i\cap U_j\) is itself affine; this is the
  canonical overlap selected by the Lean atlas.  Write
  \[
    A_i=\Gamma(U_i,\mathcal O),\qquad
    B_{ij}=\Gamma(U_{ij},\mathcal O).
  \]
  The package index is exactly
  \[
    \operatorname{Pic0FiniteStageRingIndex}
      = I\;\sqcup\;(I\times I),
  \]
  with ring \(A_i\) on the left summand and \(B_{ij}\) on the right.
  Every ring is finitely presented over \(\Omega\), and both restriction
  maps \(A_i\to B_{ij}\) and \(A_j\to B_{ij}\) are finitely presented.
  There is one finite subextension over which this complete object family is
  modeled.  Restriction and transition maps, and their equations, are the
  separate next package, not outputs of this chart-ring attachment.
\end{theorem}

\begin{proof}
  Quasi-compactness gives the finite affine cover, and separatedness of the
  represented group makes each ordered intersection affine (hence
  quasi-compact).  Local finite presentation gives finite presentation of the
  selected coordinate rings and restriction maps.  Stacks Tag 01ZM(1) descends
  this finite family.  The attached Lean producer realizes it through
  `Pic0FiniteStageChartRing`, `Pic0FiniteStageOverlapRing`, and the sum index
  displayed above.  Restriction maps, inverse transitions, and their equations
  are the separate next package, not outputs of this attachment.
\end{proof}

\begin{theorem}[The finite-stage GluePackage]
  \label{route-thm-glue-package}
  \lean{AlgebraicGeometry.Pic0FiniteStageGluePackage,
    AlgebraicGeometry.exists_pic0FiniteStageGluePackage,
    AlgebraicGeometry.Pic0FiniteStageGluePackage.presentation,
    AlgebraicGeometry.Pic0FiniteStageGluePackage.glueData,
    AlgebraicGeometry.Pic0FiniteStageGluePackage.gluedOver}
  \leanok
  \source{stacks-project}
  \dcref{Tags 01ZM, 0EUU, 081E}
  \uses{route-thm-finite-atlas}

  There is a package \(P_0\) with nested finite subextensions
  \[
    K\subseteq P_0.L.1\subseteq P_0.M.1\subseteq
      E_0:=P_0.N.1\subseteq\Omega .
  \]
  It contains finite-presentation models of every \(A_i\) and canonical
  \(B_{ij}\), all algebra maps in the restriction and reversed-overlap
  transition families, proofs that restriction maps remain open immersions,
  and the descended triple-overlap equations.  These data form
  \[
    P_0.\mathrm{glueData}:\mathrm{Scheme.GlueData}
  \]
  and a glued structure morphism
  \[
    P_0.\mathrm{gluedMap}:P_0.\mathrm{glueData.glued}
      \longrightarrow\Spec E_0.
  \]
  Set
  \(J_{E_0}=P_0.\mathrm{gluedOver}:=\Over.mk(P_0.\mathrm{gluedMap})\).
\end{theorem}

\begin{proof}
  Use Theorem~\ref{route-thm-finite-atlas} for the finite list of chart and
  canonical-overlap rings.  The `Pic0FiniteStageRestrictionIndex` supplies the
  two restriction legs for every ordered pair, while the transition and triple
  indices supply the reversal and cocycle fields.
  Enlarge to \(P_0.M.1\) so every restriction and transition map descends and
  every required open-immersion property holds; Stacks Tag 0EUU is the
  abstract finite-stage open-immersion principle.  Enlarge to \(P_0.N.1\) so
  inverse and triple-cocycle equations hold; Tag 081E is the corresponding
  isomorphism descent principle.  The contravariant spectra of the descended
  rings and maps satisfy the exact \(\mathrm{Scheme.GlueData}\) fields, whose
  standard gluing construction gives \(J_{E_0}\).
\end{proof}

\section{Nested pullbacks and overlap comparison}

\begin{theorem}[Finite-stage overlap diagrams recover the original atlas]
  \label{route-thm-overlap-comparison}
  \lean{AlgebraicGeometry.nestedPullbackFlatteningIso,
    AlgebraicGeometry.Pic0FiniteStageGluePackage.gluingOverlapFlatteningIso,
    AlgebraicGeometry.Pic0FiniteStageGluePackage.baseChangeGluingIso,
    AlgebraicGeometry.Pic0FiniteStageGluePackage.gluingOverlapIso,
    AlgebraicGeometry.Pic0FiniteStageGluePackage.gluingOverlapIso_fst,
    AlgebraicGeometry.Pic0FiniteStageGluePackage.gluingOverlapIso_snd}
  \leanok
  \dcref{custom}
  \uses{route-thm-glue-package}

  For every pair \(i,j\), scalar extension of the finite-stage overlap to
  \(\Omega\) is canonically isomorphic to the corresponding overlap
  \(U_i\cap U_j\) in \(J_\Omega\).  The isomorphism is the composite of:
  the associativity/flattening isomorphism for the nested tensor products,
  the comparison with the selected finite affine overlap cover, and the
  pullback isomorphism induced by the descended restriction maps.  Under it,
  both projections agree exactly with the left and right atlas restriction
  maps.  The transition maps agree on triple overlaps.
\end{theorem}

\begin{proof}
  On coordinate rings, associativity identifies
  \(\Omega\otimes_{E_0}
    (E_0\otimes_{P_0.M.1}(P_0.M.1\otimes_{P_0.L.1}A))\) with the original
  \(\Omega\)-algebra.  The package comparison equations identify both
  restriction homomorphisms and the cyclic transition homomorphism after this
  flattening.  Passing contravariantly to spectra gives the nested-pullback
  and overlap isomorphisms.  The named first- and second-projection theorems,
  followed by the triple face equation, verify the complete gluing diagram;
  no appeal to Picard representability occurs in this argument.
\end{proof}

\begin{theorem}[Global finite-stage base-change comparison]
  \label{route-thm-global-comparison}
  \lean{AlgebraicGeometry.Pic0FiniteStageGluePackage.gluingGluedIso,
    AlgebraicGeometry.Pic0FiniteStageGluePackage.finiteStageBaseChangeIso,
    AlgebraicGeometry.Pic0FiniteStageGluePackage.finiteStageBaseChangeIso_hom_structureMap}
  \leanok
  \source{stacks-project}
  \dcref{Tags 01JB--01JC, gluing and its mapping property; Tag 01ZM,
    finite-stage gluing}
  \uses{route-thm-overlap-comparison, route-thm-sepclosed-representability}

  The chart comparisons glue to an isomorphism over \(\Omega\)
  \[
    \beta_{P_0}:
      J_{E_0}\times_{\Spec E_0}\Spec\Omega
        \xrightarrow{\sim}J_\Omega .
  \]
  It commutes with the two structure maps to \(\Spec\Omega\).  Transporting
  \(r_\Omega\) along \(\beta_{P_0}\) proves only
  \[
    \operatorname{RepresentableBy}
      \bigl(\operatorname{Pic}^0_{C_\Omega},
            J_{E_0}\times_{E_0}\Omega\bigr).
  \]
  It does \textbf{not} prove
  \(\operatorname{RepresentableBy}
      (\operatorname{Pic}^0_{C_{E_0}},J_{E_0})\).
\end{theorem}

\begin{proof}
  Theorem~\ref{route-thm-overlap-comparison} identifies the complete
  base-changed gluing diagram with the finite atlas of \(J_\Omega\).
  Uniqueness of gluing yields the first global isomorphism; the checked
  structure-map equation upgrades it to the slice category.  The last
  displayed post-base-change certificate is ordinary transport of an
  existing representer across that slice isomorphism.  Since no Picard class
  on \(J_{E_0}\) and no natural equivalence over \(E_0\) has yet been constructed,
  interpreting this result as finite-stage representability would be a
  logical error.

  The source tree also contains the wrappers
  \texttt{finiteStageBaseChangeOverIso} and
  \texttt{pic0RepresentableBy\_finiteStageBaseChange}.  Their defining module
  is not reachable from the public \texttt{AlgebraicJacobian} root, so they are
  unverified implementation artifacts here and receive no
  \texttt{\textbackslash leanok} credit.  Their targets in any case remain the
  slice comparison and post-base-change representation over \(\Omega\).
\end{proof}
